\documentclass[11pt,a4paper]{article}

% ── Packages ──────────────────────────────────────────────────────────────────
\usepackage[utf8]{inputenc}
\usepackage[T1]{fontenc}
\usepackage[margin=2.5cm]{geometry}
\usepackage{hyperref}
\usepackage{xcolor}
\usepackage{listings}
\usepackage{enumitem}
\usepackage{booktabs}
\usepackage{longtable}
\usepackage{graphicx}
\usepackage{amsmath,amssymb}
\usepackage{fancyhdr}
\usepackage{titlesec}

% ── Colors ────────────────────────────────────────────────────────────────────
\definecolor{isls-blue}{HTML}{1a5276}
\definecolor{isls-orange}{HTML}{e67e22}
\definecolor{isls-green}{HTML}{27ae60}
\definecolor{isls-red}{HTML}{c0392b}
\definecolor{isls-gray}{HTML}{7f8c8d}
\definecolor{code-bg}{HTML}{f7f9fb}

% ── Listings ──────────────────────────────────────────────────────────────────
\lstdefinestyle{rust}{
  language=C,
  basicstyle=\ttfamily\small,
  keywordstyle=\color{isls-blue}\bfseries,
  commentstyle=\color{isls-gray}\itshape,
  stringstyle=\color{isls-green},
  backgroundcolor=\color{code-bg},
  frame=single,
  rulecolor=\color{isls-gray!30},
  breaklines=true,
  showstringspaces=false,
  tabsize=4,
  morekeywords={fn,let,mut,pub,struct,impl,enum,use,mod,self,Self,
    crate,super,where,trait,for,in,if,else,match,return,async,await,
    Ok,Err,Some,None,Vec,String,HashMap,HashSet,PathBuf,Result,Option},
}
\lstdefinestyle{bash}{
  language=bash,
  basicstyle=\ttfamily\small,
  backgroundcolor=\color{code-bg},
  frame=single,
  rulecolor=\color{isls-gray!30},
  breaklines=true,
  showstringspaces=false,
}
\lstdefinestyle{toml}{
  basicstyle=\ttfamily\small,
  backgroundcolor=\color{code-bg},
  frame=single,
  rulecolor=\color{isls-gray!30},
  breaklines=true,
  showstringspaces=false,
  commentstyle=\color{isls-gray}\itshape,
  morecomment=[l]{\#},
}

% ── Headers ───────────────────────────────────────────────────────────────────
\pagestyle{fancy}
\fancyhf{}
\fancyhead[L]{\small\textcolor{isls-gray}{ISLS D5 Spec --- Repository Scraping}}
\fancyhead[R]{\small\textcolor{isls-gray}{v1.0.0 --- \today}}
\fancyfoot[C]{\thepage}

% ── Section styling ───────────────────────────────────────────────────────────
\titleformat{\section}
  {\Large\bfseries\color{isls-blue}}{\thesection}{1em}{}
\titleformat{\subsection}
  {\large\bfseries\color{isls-blue!80}}{\thesubsection}{1em}{}
\titleformat{\subsubsection}
  {\normalsize\bfseries\color{isls-blue!60}}{\thesubsubsection}{1em}{}

\hypersetup{
  colorlinks=true,
  linkcolor=isls-blue,
  urlcolor=isls-orange,
  citecolor=isls-green,
}

% ══════════════════════════════════════════════════════════════════════════════
\begin{document}

% ── Titlepage ─────────────────────────────────────────────────────────────────
\begin{titlepage}
\centering
\vspace*{3cm}

{\Huge\bfseries\color{isls-blue} ISLS D5 Specification}\\[0.5cm]
{\LARGE\color{isls-orange} Repository Scraping: Repos $\to$ Topologie $\to$ Normen}\\[2cm]

{\large
\begin{tabular}{rl}
\textbf{System:}    & ISLS --- Intelligent Semantic Ledger Substrate \\
\textbf{Version:}   & v3.6 (D5) \\
\textbf{Author:}    & Sebastian Klemm \\
\textbf{Date:}      & \today \\
\textbf{Status:}    & Specification \\
\textbf{Depends:}   & D4 (complete), isls-reader v3.0, isls-code-topo v3.0 \\
\textbf{Branch:}    & \texttt{claude/implement-isls-d5-spec} \\
\end{tabular}
}

\vfill

{\small\color{isls-gray}
Mithras Substructure University\\
Repository: \texttt{https://github.com/LashSesh/graphity}\\
License: MIT
}
\end{titlepage}

% ── TOC ───────────────────────────────────────────────────────────────────────
\tableofcontents
\newpage

% ══════════════════════════════════════════════════════════════════════════════
\section{Executive Summary}
% ══════════════════════════════════════════════════════════════════════════════

D4 closed the feedback loop: ISLS observes its own generations and promotes
recurring patterns to norms. But the observation surface is limited to
self-generated apps --- the system only learns from itself.

D5 \textbf{extends the observation surface to the wild}. The system reads
external repositories, extracts their architectural topology using Barbara
(isls-reader) and isls-code-topo, converts the topology into
\texttt{ObservedArtifact}s, and feeds them into the same
\texttt{observe\_and\_learn()} pipeline that D4 established.

In biological terms: D4 was endogenous learning (observing own metabolism).
D5 is exogenous feeding --- the mycelium extends hyphae into foreign
substrates, digests the structural topology, and incorporates the nutrients
into its own norm catalog.

\subsection{The Legal Boundary}

D5 extracts \textbf{topology}, not code:

\begin{itemize}[nosep]
\item Function signatures: ``there is a public async function named
      \texttt{get\_product} with 2 parameters'' --- not the function body.
\item Struct names and field names: ``there is a struct \texttt{Product}
      with fields \texttt{name}, \texttt{price}, \texttt{quantity}'' ---
      not the implementations.
\item Layer organization: ``the project has models/, services/, api/
      directories'' --- not the file contents.
\item Import graph: ``file A imports module B'' --- not what module B does.
\end{itemize}

Architectural patterns are not copyrightable. ISLS never stores, copies,
or reproduces source code from scraped repositories. The
\texttt{ObservedArtifact} contains a SHA-256 content hash for deduplication
--- not the content itself.

\subsection{What D5 Achieves}

\begin{enumerate}[nosep]
\item \textbf{Topology-to-Observation Bridge:} Converts
      \texttt{CodeObservation} (Barbara output) into
      \texttt{ObservedArtifact} (norm learning input).
\item \textbf{Local Scraping:} \texttt{isls scrape --path ./some-repo}
      reads a local directory.
\item \textbf{Git Scraping:} \texttt{isls scrape --url https://github.com/...}
      shallow-clones and scrapes.
\item \textbf{Manifest Scraping:} \texttt{isls scrape --manifest repos.toml}
      batch-processes a list of repositories.
\item \textbf{Scrape Report:} Summary of extracted topology, entities
      found, artifacts generated, candidates updated.
\end{enumerate}

\subsection{What D5 is NOT}

\begin{itemize}[nosep]
\item \textbf{Not code generation.} D5 only observes --- it does not
      generate applications.
\item \textbf{Not code copying.} No source code is stored or reproduced.
\item \textbf{Not LLM-dependent.} D5 uses zero LLM tokens. Pure structural
      analysis.
\item \textbf{Not a package manager.} D5 does not install dependencies or
      build scraped projects.
\item \textbf{Not self-generation.} That is D6.
\end{itemize}

\subsection{Token Budget}

D5 costs \textbf{zero API tokens}. All analysis is structural (regex
parsing via Barbara, topology computation, pattern extraction). This
makes D5 unlimited in scale --- the only constraint is disk space for
temporary clones and CPU time for parsing.


% ══════════════════════════════════════════════════════════════════════════════
\section{Architecture Overview}
% ══════════════════════════════════════════════════════════════════════════════

\subsection{Existing Infrastructure}

Two crates already exist but are underutilized:

\paragraph{isls-reader (Barbara)} Multi-language source code parser.
Produces \texttt{CodeObservation} containing functions, structs, imports,
SQL tables, LOC, and SHA-256 per file. Supports Rust, JavaScript/TypeScript,
Python, SQL, HTML, CSS, TOML, Dockerfile.

Key APIs:
\begin{itemize}[nosep]
\item \texttt{parse\_file(path)} --- single file $\to$ \texttt{CodeObservation}
\item \texttt{parse\_directory(dir)} --- recursive $\to$ \texttt{WorkspaceAnalysis}
\item \texttt{Language::from\_path()} --- auto-detection from extension
\end{itemize}

\paragraph{isls-code-topo} Computes topological signatures from
\texttt{CodeObservation} sets. Produces \texttt{CodeTopology} containing
function signatures, struct names, layer detection, import graph, and
connectivity measure. Also provides \texttt{topology\_similarity()} for
comparing two topologies (Jaccard-based, hardened weights from Barbara).

\subsection{The Missing Bridge}

Currently there is no path from:
\begin{center}
\texttt{CodeObservation} (Barbara) $\to$ \texttt{ObservedArtifact} (norm learning)
\end{center}

This bridge is the core of D5. Once it exists, any codebase that Barbara
can read becomes food for the norm system.

\subsection{Target Data Flow}

\begin{lstlisting}[style=bash]
Repository (local or git URL)
    |
    v
isls-reader::parse_directory()  -->  Vec<CodeObservation>
    |
    v
isls-code-topo::compute_code_topology()  -->  CodeTopology
    |                                           (for reporting)
    v
topo_bridge::observations_from_code()  -->  Vec<ObservedArtifact>
    |
    v
NormRegistry::observe_and_learn()  -->  Candidates updated
    |
    v
~/.isls/norms.json persisted
\end{lstlisting}


% ══════════════════════════════════════════════════════════════════════════════
\section{Work Packages}
% ══════════════════════════════════════════════════════════════════════════════

% ──────────────────────────────────────────────────────────────────────────────
\subsection{W1: Topology-to-Observation Bridge}
\label{sec:w1}
% ──────────────────────────────────────────────────────────────────────────────

\subsubsection{Goal}

Create a module that converts Barbara's \texttt{CodeObservation}s into
\texttt{ObservedArtifact}s compatible with D4's \texttt{observe\_and\_learn()}.

\subsubsection{New File}

\begin{longtable}{lp{10cm}}
\toprule
\textbf{File} & \textbf{Purpose} \\
\midrule
\texttt{isls-code-topo/src/bridge.rs}
  & Converts \texttt{Vec<CodeObservation>} into
    \texttt{Vec<ObservedArtifact>}. Pure function, no IO. \\
\bottomrule
\end{longtable}

\subsubsection{Conversion Algorithm}

For each \texttt{CodeObservation}:

\begin{enumerate}[nosep]
\item \textbf{Determine LayerType} from file path using the existing
      \texttt{detect\_layer()} logic in isls-code-topo, extended:
      \begin{itemize}[nosep]
      \item \texttt{models/}, \texttt{model/}, \texttt{entities/}
            $\to$ \texttt{LayerType::Model}
      \item \texttt{database/}, \texttt{db/}, \texttt{queries/},
            \texttt{repository/}, \texttt{repositories/}
            $\to$ \texttt{LayerType::Query}
      \item \texttt{services/}, \texttt{service/}, \texttt{domain/}
            $\to$ \texttt{LayerType::Service}
      \item \texttt{api/}, \texttt{routes/}, \texttt{handlers/},
            \texttt{controllers/}, \texttt{endpoints/}
            $\to$ \texttt{LayerType::Api}
      \item \texttt{frontend/}, \texttt{pages/}, \texttt{components/},
            \texttt{views/}, \texttt{templates/}
            $\to$ \texttt{LayerType::Frontend}
      \item \texttt{tests/}, \texttt{test/}, \texttt{\_test.},
            \texttt{\_tests.}, \texttt{spec/}
            $\to$ \texttt{LayerType::Test}
      \item \texttt{migrations/}, \texttt{schema.sql}, \texttt{*.sql}
            in root $\to$ \texttt{LayerType::Database}
      \item \texttt{config/}, \texttt{settings/},
            \texttt{*.toml}/\texttt{*.yaml} in root
            $\to$ \texttt{LayerType::Config}
      \item Everything else $\to$ skip (no observation generated)
      \end{itemize}

\item \textbf{Extract ObservedArtifacts} from the observation:
      \begin{itemize}[nosep]
      \item For each \texttt{StructDef}: one \texttt{ObservedArtifact}
            with \texttt{artifact\_type = "struct"},
            \texttt{name = struct\_name},
            \texttt{field\_names = struct.fields},
            \texttt{signature = file SHA-256}.
      \item For each \texttt{FunctionDef}: one \texttt{ObservedArtifact}
            with \texttt{artifact\_type = "fn"},
            \texttt{name = function\_name},
            \texttt{field\_names = function.params},
            \texttt{signature = file SHA-256}.
      \item For each \texttt{SqlTable}: one \texttt{ObservedArtifact}
            with \texttt{layer = Database},
            \texttt{artifact\_type = "table"},
            \texttt{name = table\_name},
            \texttt{field\_names = table.columns}.
      \end{itemize}
\end{enumerate}

\subsubsection{Domain Name Inference}

The ``domain'' passed to \texttt{observe\_and\_learn()} is inferred from
the repository:

\begin{itemize}[nosep]
\item Local path: last directory component
      (e.g. \texttt{./repos/petshop-api} $\to$ \texttt{"petshop-api"})
\item Git URL: repository name from URL
      (e.g. \texttt{https://github.com/user/petshop-api.git}
      $\to$ \texttt{"petshop-api"})
\end{itemize}

\subsubsection{Public API}

\begin{lstlisting}[style=rust]
/// Convert Barbara CodeObservations into ObservedArtifacts
/// for the norm learning pipeline.
///
/// Skips files that cannot be mapped to a known LayerType.
/// Returns (artifacts, skipped_count).
pub fn observations_from_code(
    code_obs: &[CodeObservation],
) -> (Vec<ObservedArtifact>, usize)
\end{lstlisting}

\subsubsection{Modified Files}

\begin{longtable}{lp{10cm}}
\toprule
\textbf{File} & \textbf{Change} \\
\midrule
\texttt{isls-code-topo/src/lib.rs}
  & Add \texttt{pub mod bridge;} \\
\texttt{isls-code-topo/Cargo.toml}
  & Add dependency: \texttt{isls-norms = \{path = "../isls-norms"\}}
    (for \texttt{ObservedArtifact} and \texttt{LayerType}) \\
\bottomrule
\end{longtable}

\subsubsection{W1 Verification}

\begin{enumerate}[nosep]
\item Unit test: parse a mock Rust project (3 files: model, query,
      service) via \texttt{parse\_string()}, convert with bridge.
\item Assert correct LayerType assignment for each artifact.
\item Assert struct fields appear in \texttt{field\_names}.
\item Assert function params appear in \texttt{field\_names}.
\item Assert files in unmapped directories produce no artifacts
      (skipped\_count $> 0$).
\end{enumerate}


% ──────────────────────────────────────────────────────────────────────────────
\subsection{W2: Local Scraping}
\label{sec:w2}
% ──────────────────────────────────────────────────────────────────────────────

\subsubsection{Goal}

Add the \texttt{isls scrape} CLI command for scraping local directories.

\subsubsection{Command}

\begin{lstlisting}[style=bash]
isls scrape --path ./path/to/repo
isls scrape --path ./path/to/repo --domain "custom-domain-name"
\end{lstlisting}

\subsubsection{Pipeline}

\begin{enumerate}[nosep]
\item Validate path exists and is a directory.
\item Call \texttt{isls\_reader::parse\_directory(path)} $\to$
      \texttt{WorkspaceAnalysis}.
\item Call \texttt{bridge::observations\_from\_code(\&analysis.files)}
      $\to$ \texttt{Vec<ObservedArtifact>}.
\item Compute \texttt{CodeTopology} for the report.
\item Infer domain name from path (or use \texttt{--domain} override).
\item Generate run\_id: \texttt{scrape\_\{domain\}\_\{timestamp\}}.
\item Load \texttt{NormRegistry}, call
      \texttt{observe\_and\_learn(\&artifacts, domain, run\_id)}.
\item Print scrape report.
\end{enumerate}

\subsubsection{Scrape Report Format}

\begin{lstlisting}[style=bash]
ISLS Scrape Report
---
Source:         ./repos/petshop-api
Domain:         petshop-api
Languages:      rust (1842 LOC), sql (156 LOC)
Files parsed:   23
Artifacts:      47 (12 struct, 28 fn, 7 table)
Skipped:        6 (unmapped directories)
Layers:         Model(8), Query(12), Service(8), Api(11), Database(7), Test(1)
---
Topology:
  Nodes: 47, Edges: 31, Connectivity: 0.014
  Structs: Product, Order, Category, User, ...
  Functions: get_product/2, list_products/1, create_product/3, ...
---
Norm System:
  Candidates updated: 3
  New candidates: 1
  Promoted: 0
\end{lstlisting}

\subsubsection{New Files}

\begin{longtable}{lp{10cm}}
\toprule
\textbf{File} & \textbf{Purpose} \\
\midrule
\texttt{isls-cli/src/cmd\_scrape.rs}
  & Implements \texttt{cmd\_scrape()} --- local and git scraping,
    manifest batch processing, report generation. \\
\bottomrule
\end{longtable}

\subsubsection{Modified Files}

\begin{longtable}{lp{10cm}}
\toprule
\textbf{File} & \textbf{Change} \\
\midrule
\texttt{isls-cli/src/main.rs}
  & Add \texttt{Command::Scrape} variant. Parse \texttt{isls scrape}
    with \texttt{--path}, \texttt{--url}, \texttt{--manifest},
    \texttt{--domain} options. \\
\texttt{isls-cli/Cargo.toml}
  & Add dependencies: \texttt{isls-reader}, \texttt{isls-code-topo} \\
\bottomrule
\end{longtable}

\subsubsection{W2 Verification}

\begin{enumerate}[nosep]
\item Scrape one of the previously generated ISLS output directories:
      \texttt{isls scrape --path ./output/pipeline\_*/petshop}
\item Verify report shows correct language, file count, artifact count.
\item Verify \texttt{isls norms candidates} shows updated observation
      counts (the petshop domain now has observations from both
      D4 generation AND D5 scraping).
\item Scrape the ISLS source itself:
      \texttt{isls scrape --path ./isls --domain "isls-self"}
\item Verify report shows Rust as primary language with correct
      struct/function counts.
\end{enumerate}


% ──────────────────────────────────────────────────────────────────────────────
\subsection{W3: Git Clone Scraping}
\label{sec:w3}
% ──────────────────────────────────────────────────────────────────────────────

\subsubsection{Goal}

Support scraping repositories directly from Git URLs. The clone is
temporary, shallow, and automatically cleaned up.

\subsubsection{Command}

\begin{lstlisting}[style=bash]
isls scrape --url https://github.com/user/repo.git
isls scrape --url https://github.com/user/repo.git --domain "custom-name"
\end{lstlisting}

\subsubsection{Clone Strategy}

\begin{enumerate}[nosep]
\item Create temp directory:
      \texttt{std::env::temp\_dir().join("isls-scrape-\{hash\}")}.
\item Execute: \texttt{git clone --depth 1 --single-branch <url> <temp>}.
\item Verify clone succeeded (check exit code).
\item Proceed with local scraping pipeline (W2) on temp directory.
\item Delete temp directory after scraping (cleanup in all paths,
      including errors --- use RAII/drop guard).
\end{enumerate}

\subsubsection{Constraints for i3 ThinkPad}

\begin{itemize}[nosep]
\item \textbf{Shallow clone only} (\texttt{--depth 1}). No history.
\item \textbf{Size limit:} If the cloned directory exceeds 200 MB,
      abort with a warning (not an error --- the repo is just too big
      for this machine). Configurable via \texttt{--max-size-mb}.
\item \textbf{Timeout:} Clone must complete within 120 seconds.
      Configurable via \texttt{--timeout}.
\item \textbf{Cleanup guarantee:} Temp directory MUST be deleted on all
      paths (success, error, panic). Use a guard struct with \texttt{Drop}.
\end{itemize}

\subsubsection{Git Availability Check}

Before cloning, verify that \texttt{git} is available:
\texttt{std::process::Command::new("git").arg("--version").status()}.
If not found, print a clear error message and exit.

\subsubsection{W3 Verification}

\begin{enumerate}[nosep]
\item \texttt{isls scrape --url https://github.com/LashSesh/graphity.git
      --domain "graphity"}
\item Verify clone succeeds, scrape report shows multi-language analysis.
\item Verify temp directory is cleaned up (does not exist after command
      completes).
\item Test with a non-existent URL: verify graceful error message.
\item Test with \texttt{--max-size-mb 1}: verify abort on large repos.
\end{enumerate}


% ──────────────────────────────────────────────────────────────────────────────
\subsection{W4: Manifest Batch Scraping}
\label{sec:w4}
% ──────────────────────────────────────────────────────────────────────────────

\subsubsection{Goal}

Support batch scraping from a TOML manifest file listing multiple
repositories.

\subsubsection{Command}

\begin{lstlisting}[style=bash]
isls scrape --manifest repos.toml
\end{lstlisting}

\subsubsection{Manifest Format}

\begin{lstlisting}[style=toml]
# repos.toml — ISLS Scrape Manifest
# Each [[repo]] entry is scraped sequentially.

[[repo]]
url = "https://github.com/user/project-a.git"
domain = "project-a"

[[repo]]
url = "https://github.com/user/project-b.git"
# domain auto-inferred from URL

[[repo]]
path = "./local-repos/project-c"
domain = "project-c"

[[repo]]
url = "https://github.com/user/large-project.git"
max_size_mb = 500   # override default 200 MB limit
\end{lstlisting}

\subsubsection{Batch Pipeline}

\begin{enumerate}[nosep]
\item Parse manifest TOML.
\item For each \texttt{[[repo]]} entry:
      \begin{enumerate}[nosep,label=\alph*.]
      \item Display progress: \texttt{[3/12] Scraping project-a...}
      \item Execute local scrape (if \texttt{path}) or git scrape
            (if \texttt{url}).
      \item On failure: log warning, continue to next repo.
      \item Accumulate per-repo results.
      \end{enumerate}
\item After all repos: print batch summary report.
\item Call \texttt{isls norms stats} equivalent at end.
\end{enumerate}

\subsubsection{Batch Report Format}

\begin{lstlisting}[style=bash]
ISLS Batch Scrape Report
---
Manifest:  repos.toml (12 repos)

  [PASS] project-a    rust    1842 LOC   47 artifacts
  [PASS] project-b    python  3201 LOC   89 artifacts
  [FAIL] project-c    clone timeout
  [PASS] project-d    js+ts   5603 LOC   134 artifacts
  ...

Summary:
  Scraped:     11/12
  Total LOC:   42,871
  Total artifacts: 891
  Languages:   rust(5), python(3), javascript(2), typescript(1)

Norm System:
  Candidates before: 5
  Candidates after:  14
  New auto-norms:    3
\end{lstlisting}

\subsubsection{Shipped Manifest}

Include a default manifest in the repository:

\begin{lstlisting}[style=bash]
isls/examples/scrape_manifest.toml
\end{lstlisting}

Content: 5--10 well-known open-source full-stack Rust projects
(e.g. Actix-Web examples, Axum examples, Rocket examples). These
serve as seed data for the norm catalog.

\subsubsection{W4 Verification}

\begin{enumerate}[nosep]
\item Create a test manifest with 3 entries (2 local paths from
      D4 output + 1 git URL).
\item \texttt{isls scrape --manifest test\_manifest.toml}
\item Verify batch report shows 2--3 PASS entries.
\item Verify norm candidates accumulated across all repos.
\end{enumerate}


% ══════════════════════════════════════════════════════════════════════════════
\section{File Inventory}
% ══════════════════════════════════════════════════════════════════════════════

\subsection{New Files}

\begin{longtable}{llp{7cm}}
\toprule
\textbf{Crate} & \textbf{File} & \textbf{Purpose} \\
\midrule
isls-code-topo & \texttt{src/bridge.rs}
  & Topology-to-Observation converter \\
isls-cli & \texttt{src/cmd\_scrape.rs}
  & Scrape command: local, git, manifest \\
isls & \texttt{examples/scrape\_manifest.toml}
  & Default scrape manifest (seed repos) \\
\bottomrule
\end{longtable}

\subsection{Modified Files}

\begin{longtable}{llp{7cm}}
\toprule
\textbf{Crate} & \textbf{File} & \textbf{Change} \\
\midrule
isls-code-topo & \texttt{src/lib.rs}
  & Add \texttt{pub mod bridge;} \\
isls-code-topo & \texttt{Cargo.toml}
  & Add isls-norms dependency \\
isls-cli & \texttt{src/main.rs}
  & Add Scrape command variant + argument parsing \\
isls-cli & \texttt{Cargo.toml}
  & Add isls-reader, isls-code-topo dependencies \\
\bottomrule
\end{longtable}

\subsection{Unchanged}

All of isls-forge-llm (HDAG, structural, oracle, staged\_closure),
isls-norms (types, catalog, composition, learning),
isls-chat, isls-reader, isls-hypercube, isls-types
remain \textbf{unchanged}. D5 is a pure consumer of existing APIs.


% ══════════════════════════════════════════════════════════════════════════════
\section{The 12 Rules --- D5 Compliance}
% ══════════════════════════════════════════════════════════════════════════════

\begin{longtable}{clp{9cm}}
\toprule
\textbf{\#} & \textbf{Status} & \textbf{D5 Impact} \\
\midrule
1  & \color{isls-green}\checkmark & D5 does not generate code.
     No structural/LLM boundary affected. \\
2  & \color{isls-green}\checkmark & New module \texttt{bridge.rs}
     re-exported via \texttt{pub mod bridge;} in lib.rs. \\
3  & \color{isls-green}\checkmark & Extracted function names preserve
     original naming. No renaming. \\
4  & \color{isls-green}\checkmark & No auth changes. \\
5  & \color{isls-green}\checkmark & No query changes. \\
6  & \color{isls-green}\checkmark & No query file changes. \\
7  & \color{isls-green}\checkmark & No new crypto or logging crates. \\
8  & \color{isls-green}\checkmark & No seed changes. \\
9  & \color{isls-green}\checkmark & No Dockerfile changes. \\
10 & \color{isls-green}\checkmark & Clone failure, parse failure,
     oversized repo $\to$ warning + continue (graceful degradation). \\
11 & \color{isls-green}\checkmark & No ProvidedSymbol changes. \\
12 & \color{isls-green}\checkmark & No LLM prompt changes. D5 uses
     zero LLM tokens. \\
\bottomrule
\end{longtable}


% ══════════════════════════════════════════════════════════════════════════════
\section{Data Flow: Complete D5 Cycle}
% ══════════════════════════════════════════════════════════════════════════════

\begin{enumerate}
\item \textbf{User:}
      \texttt{isls scrape --url https://github.com/example/todo-api.git}

\item \textbf{isls-cli (Clone):}
      \begin{enumerate}[nosep,label=\alph*.]
      \item Verify git is available.
      \item Create temp dir \texttt{/tmp/isls-scrape-a7f3c2/}.
      \item \texttt{git clone --depth 1 --single-branch <url> <temp>}.
      \item Check size $<$ 200 MB.
      \end{enumerate}

\item \textbf{isls-reader (Barbara):}
      \begin{enumerate}[nosep,label=\alph*.]
      \item \texttt{parse\_directory(temp\_dir)} $\to$ WorkspaceAnalysis.
      \item Result: 15 CodeObservations (8 Rust, 3 SQL, 2 JS, 2 TOML).
      \end{enumerate}

\item \textbf{isls-code-topo (Bridge):}
      \begin{enumerate}[nosep,label=\alph*.]
      \item \texttt{observations\_from\_code(\&analysis.files)} $\to$
            31 ObservedArtifacts.
      \item 4 files skipped (unmapped directories).
      \item Also: \texttt{compute\_code\_topology()} for the report.
      \end{enumerate}

\item \textbf{isls-norms (Observation):}
      \begin{enumerate}[nosep,label=\alph*.]
      \item Load NormRegistry from \texttt{\textasciitilde/.isls/norms.json}.
      \item \texttt{observe\_and\_learn(\&artifacts, "todo-api",
            "scrape\_todo-api\_20260402\_1230")}.
      \item Cross-layer patterns extracted, candidates updated.
      \item Save to disk.
      \end{enumerate}

\item \textbf{Cleanup:}
      Delete \texttt{/tmp/isls-scrape-a7f3c2/}.

\item \textbf{Report:}
      Print scrape summary showing languages, artifacts, layers,
      norm system status.
\end{enumerate}


% ══════════════════════════════════════════════════════════════════════════════
\section{Guard Struct for Temp Directory Cleanup}
% ══════════════════════════════════════════════════════════════════════════════

The temp directory MUST be cleaned up in all cases. Use an RAII guard:

\begin{lstlisting}[style=rust]
struct TempCloneGuard {
    path: PathBuf,
}

impl Drop for TempCloneGuard {
    fn drop(&mut self) {
        if self.path.exists() {
            let _ = std::fs::remove_dir_all(&self.path);
        }
    }
}
\end{lstlisting}

Create the guard immediately after creating the temp directory.
When the guard goes out of scope (including on panic or early return),
the directory is deleted.


% ══════════════════════════════════════════════════════════════════════════════
\section{Error Handling}
% ══════════════════════════════════════════════════════════════════════════════

\begin{longtable}{lp{10cm}}
\toprule
\textbf{Failure Mode} & \textbf{Behavior} \\
\midrule
Path does not exist & Print error, exit with code 1. \\
Path is a file, not dir & Print error, suggest correct usage. \\
Git not installed & Print ``git is required for URL scraping.
  Install git or use --path for local directories.'' \\
Clone fails (network) & Print error, suggest checking URL. In
  manifest mode: log warning, continue to next repo. \\
Clone timeout (120s) & Kill process, print timeout warning.
  Temp dir cleaned up via guard. \\
Repo exceeds size limit & Print warning with actual size.
  Temp dir cleaned up. \\
Barbara parse errors & Individual file errors logged, other
  files still processed. \\
No artifacts extracted & Print warning ``No architectural patterns
  found. The project may use non-standard directory layout.'' \\
\texttt{norms.json} corrupt & Handled by NormRegistry (D4):
  fresh start with builtins. \\
\bottomrule
\end{longtable}


% ══════════════════════════════════════════════════════════════════════════════
\section{Test Suite}
% ══════════════════════════════════════════════════════════════════════════════

\subsection{Unit Tests}

\begin{longtable}{lp{10cm}}
\toprule
\textbf{Test} & \textbf{Assertion} \\
\midrule
\texttt{test\_bridge\_rust\_model}
  & Rust struct in \texttt{models/} $\to$ ObservedArtifact with
    LayerType::Model, artifact\_type ``struct'', correct fields. \\
\texttt{test\_bridge\_rust\_query}
  & Rust fn in \texttt{database/} $\to$ LayerType::Query,
    artifact\_type ``fn'', correct params. \\
\texttt{test\_bridge\_sql\_table}
  & SQL CREATE TABLE $\to$ LayerType::Database,
    artifact\_type ``table'', correct columns. \\
\texttt{test\_bridge\_python\_controller}
  & Python def in \texttt{api/} $\to$ LayerType::Api. \\
\texttt{test\_bridge\_js\_component}
  & JS class in \texttt{components/} $\to$ LayerType::Frontend. \\
\texttt{test\_bridge\_skip\_unmapped}
  & File in \texttt{utils/} $\to$ skipped, not in artifacts. \\
\texttt{test\_bridge\_multi\_entity}
  & Project with 3 entities across 4 layers $\to$ 12+ artifacts
    with correct entity grouping. \\
\texttt{test\_domain\_from\_path}
  & \texttt{./repos/my-app} $\to$ ``my-app''. \\
\texttt{test\_domain\_from\_url}
  & \texttt{https://github.com/u/my-app.git} $\to$ ``my-app''. \\
\texttt{test\_manifest\_parse}
  & Valid TOML manifest parsed to correct repo entries. \\
\bottomrule
\end{longtable}

\subsection{Integration Tests}

\begin{longtable}{lp{10cm}}
\toprule
\textbf{Test} & \textbf{Assertion} \\
\midrule
\texttt{test\_scrape\_self\_generated}
  & Scrape a D4 output directory. Verify artifacts produced,
    candidates updated. \\
\texttt{test\_scrape\_isls\_source}
  & \texttt{isls scrape --path ./isls} works without panic,
    produces report. \\
\bottomrule
\end{longtable}

\subsection{Acceptance Test (Manual)}

\begin{lstlisting}[style=bash]
# Step 1: Scrape a D4-generated app (no network needed)
isls scrape --path ./output/pipeline_*/petshop

# Verify: report shows artifacts, norms candidates updated

# Step 2: Scrape ISLS itself
isls scrape --path ./isls --domain "isls-self"

# Verify: report shows Rust, ~170+ functions, ~50+ structs

# Step 3: Scrape from GitHub (requires network + git)
isls scrape --url https://github.com/LashSesh/graphity.git

# Verify: clone, scrape, cleanup, report

# Step 4: Batch scrape from manifest
isls scrape --manifest examples/scrape_manifest.toml

# Verify: batch report, multiple repos processed

# Step 5: Verify norm accumulation
isls norms stats
isls norms candidates

# Expected: more candidates than before, observations
# from both D4 generations AND D5 scrapes
\end{lstlisting}


% ══════════════════════════════════════════════════════════════════════════════
\section{Constraints for Claude Code}
% ══════════════════════════════════════════════════════════════════════════════

\begin{enumerate}
\item \textbf{Do not modify} isls-forge-llm, isls-norms, isls-chat,
      isls-reader, isls-types, or isls-hypercube.
\item \textbf{bridge.rs} MUST be pure --- no IO, no filesystem access.
      Input: \texttt{\&[CodeObservation]}. Output:
      \texttt{(Vec<ObservedArtifact>, usize)}. All IO happens in
      \texttt{cmd\_scrape.rs}.
\item \textbf{No new external crates.} Git clone uses
      \texttt{std::process::Command}. TOML parsing uses \texttt{toml}
      (already in workspace). Temp dirs use \texttt{std::env::temp\_dir()}.
\item \textbf{No source code stored.} The bridge MUST NOT include file
      contents, function bodies, or any copyrightable text in the
      \texttt{ObservedArtifact}. Only names, signatures, types, and
      structural metadata.
\item \textbf{Graceful degradation everywhere.} Every external operation
      (clone, parse, read) must handle failure without panic.
\item \textbf{TempCloneGuard} MUST be used for git clones. No
      exception.
\item All new code MUST compile with \texttt{cargo build --workspace}.
\item Test with \texttt{cargo test -p isls-code-topo} and
      \texttt{cargo test -p isls-cli}.
\item Branch: \texttt{claude/implement-isls-d5-spec}.
\item Commit message format: \texttt{D5/W\{n\}: <description>}.
\end{enumerate}


% ══════════════════════════════════════════════════════════════════════════════
\section{Dependency Graph}
% ══════════════════════════════════════════════════════════════════════════════

\begin{verbatim}
W1 (Topology-to-Observation Bridge)
  |
  v
W2 (Local Scraping)
  |
  v
W3 (Git Clone Scraping)
  |
  v
W4 (Manifest Batch Scraping)
\end{verbatim}

W1 is the foundation --- without the bridge, no scraping can feed norms.
W2 is the simplest scraping mode (no network). W3 adds git.
W4 adds batch processing.

Recommended implementation order: W1 $\to$ W2 $\to$ W3 $\to$ W4.


% ══════════════════════════════════════════════════════════════════════════════
\section{Norm Accumulation Strategy}
% ══════════════════════════════════════════════════════════════════════════════

D5 scrapes feed into the same \texttt{observe\_and\_learn()} as D4
generations. This means:

\begin{itemize}[nosep]
\item Patterns from external repos and self-generated apps merge into
      the same candidates.
\item A pattern observed in a scraped Actix-Web project and in a
      self-generated ISLS app share the same candidate --- because
      they have the same layer topology.
\item External observations accelerate promotion: if ISLS has already
      observed a 5-layer CRUD pattern in 3 self-generated apps (D4),
      scraping 2 external repos with the same pattern pushes the
      candidate past \texttt{min\_observations: 5}.
\end{itemize}

This is the mycelial principle: endogenous and exogenous nutrients
converge in the same metabolic pathway.


% ══════════════════════════════════════════════════════════════════════════════
\section{Verification Sequence (Complete)}
\label{sec:verification}
% ══════════════════════════════════════════════════════════════════════════════

\begin{enumerate}
\item \texttt{cargo build --workspace} --- compiles.
\item \texttt{cargo test -p isls-code-topo} --- bridge tests pass.
\item \texttt{cargo test -p isls-cli} --- scrape CLI tests pass.
\item \texttt{cargo test --workspace} --- all tests pass (171 + new).
\item \texttt{isls scrape --path ./output/pipeline\_*/petshop}
      --- local scrape produces report.
\item \texttt{isls norms candidates} --- shows updated observations.
\item \texttt{isls scrape --path ./isls --domain "isls-self"}
      --- self-scrape works.
\item \texttt{isls scrape --url https://github.com/LashSesh/graphity.git}
      --- git scrape works (requires network).
\item Temp directory cleaned up after git scrape.
\item \texttt{isls scrape --manifest examples/scrape\_manifest.toml}
      --- batch works.
\item \texttt{isls norms stats} --- shows accumulated observations
      from both D4 and D5 sources.
\end{enumerate}

\vfill
\begin{center}
\color{isls-gray}\rule{0.5\textwidth}{0.4pt}\\[0.5em]
{\small End of D5 Specification.}
\end{center}

\end{document}
